%============================================================
% BAB V — PENUTUP
% Struktur: A. Kesimpulan, B. Saran.
% Teks disalin verbatim dari template asli
% (reference/Template_Skripsi.docx).
%============================================================
\chapter{PENUTUP}

\section{Kesimpulan}
\begin{isi}
Kesimpulan merupakan pernyataan singkat, jelas, dan tepat tentang
apa yang diperoleh, memuat keunggulan dan kelemahan, dapat
dibuktikan, serta terkait langsung dengan Rumusan Masalah dan
Tujuan Penelitian. Uraian pada bagian ini harus merupakan
pernyataan yang pernah dianalisis/dibahas pada bagian sebelumnya,
bukan pernyataan yang sama sekali baru dan tidak pernah dibahas
pada bagian sebelumnya, serta merupakan jawaban atas
permasalahan yang dirumuskan. Bagian ini tidak perlu ada uraian
penjelasan lagi.
\end{isi}

\section{Saran}
\begin{isi}
Saran memuat berbagai usulan atau pendapat yang sebaiknya
dikaitkan oleh penelitian sejenis. Saran dibuat berdasarkan
kelemahan, pengalaman, kesulitan, kesalahan, temuan baru yang
belum diteliti dan berbagai kemungkinan arah penelitian
selanjutnya. Saran yang disusun harus didasarkan atas hasil
penelitian yang dilakukan.
\end{isi}
